% ch4_a 第4章 §4.1–§4.2 (书页 144–154 / p159–p169)
% 第4章 FREE FERMION SYSTEMS
\chapter{自由费米子系统}

费米子系统是凝聚态物理中最重要的系统之一。金属、半导体、磁体、超导体等等都是费米子系统，它们的性质主要受电子的费米统计控制。在本章中，我们将研究自由多费米子系统的一些性质。

\section{多费米子系统}

\subsection{什么是费米子？}

\begin{keypoints}
\item 费米子由泡利不相容原理以及一种在两个全同费米子交换时产生$\pi$相移的跃迁哈密顿量来表征。
\item 费米子之所以怪异，是因为它们是非局域的客体。
\item 费米子可以用反对易算符来描述。
\end{keypoints}

长期以来，我们一直觉得自己知道费米子是什么。在读完接下来几节之后，如果你开始觉得自己并不真正理解费米子究竟是什么，那么我就达到了目的。我们将在第10章的后面对"费米子是什么"给出一个回答。这里我们只用传统的方式引进费米子，并在传统图像之内展示费米子是多么奇怪、多么不自然。

我们先尝试描述格点上的无自旋费米子系统。由于泡利不相容原理，格点费米子系统的希尔伯特空间可以用基矢$\{|n_{\pmb{i}_1},n_{\pmb{i}_2},\ldots\rangle\}$展开，其中$n_{\pmb{i}}=0,1$是格点$\pmb{i}$上的费米子数目。为了对自由费米子系统给出二次量子化的描述，我们为每个格点$\pmb{i}$引入湮灭算符$\sigma_{\pmb{i}}^{-}$和产生算符$\sigma_{\pmb{i}}^{+}$。它们具有如下矩阵形式：
\begin{equation*}
\sigma_{\pmb{i}}^{-}=\begin{pmatrix}0 & 0\\ 1 & 0\end{pmatrix}=\frac{1}{2}(\sigma^{x}-\mathrm{i}\sigma^{y})
\qquad
\sigma_{\pmb{i}}^{+}=\begin{pmatrix}0 & 1\\ 0 & 0\end{pmatrix}=\frac{1}{2}(\sigma^{x}+\mathrm{i}\sigma^{y})
\end{equation*}
其中$\sigma^{x,y,z}$是泡利矩阵。湮灭算符$\sigma^{-}$把有一个费米子的态$|1\rangle=\binom{1}{0}$变成没有费米子的态$|0\rangle=\binom{0}{1}$。由于$\sigma_{\pmb{i}}^{\pm}$产生/湮灭一个费米子，我们暂且把它们称作"费米子"算符。

利用"费米子"算符$\sigma_{\pmb{i}}^{\pm}$，我们可以为费米系统写下如下哈密顿量：
\begin{equation*}
H_b=\sum_{\langle\pmb{i}\pmb{j}\rangle}(t_{\pmb{i}\pmb{j}}\sigma_{\pmb{i}}^{+}\sigma_{\pmb{j}}^{-}+\text{h.c.})
\end{equation*}
虽然从数学上看，$H_b$是作用于费米子希尔伯特空间$\{|n_{\pmb{i}_1},n_{\pmb{i}_2},\ldots\rangle\}$之内的厄米算符，但由$H_b$所描述的系统并不是费米子系统，它其实是一个硬核玻色子（hard-core boson）系统或自旋-1/2系统。

这是因为我们的费米子希尔伯特空间同样可以看作一个硬核玻色子的希尔伯特空间，其中$|0\rangle$是零玻色子态，$|1\rangle$是单玻色子态。\footnote{费米子希尔伯特空间也可以看作一个自旋-1/2的希尔伯特空间，其中$|0\rangle$是自旋向下态，$|1\rangle$是自旋向上态。}所以，单凭希尔伯特空间本身，我们无法判断该系统是费米子系统还是玻色子系统；我们必须考察哈密顿量，才能判断系统是费米型的还是玻色型的。哈密顿量$H_b$是用$\sigma_{\pmb{i}}^{\pm}$写出的，而$\sigma_{\pmb{i}}^{\pm}$在不同格点上彼此对易，因此由$H_b$描述的系统是玻色子系统。我们其实应该把$\sigma_{\pmb{i}}^{-}$称作玻色子算符。

令人十分惊讶的是，费米子希尔伯特空间中一个自然的局域哈密顿量所描述的竟然不是费米子系统！这就引出一个有趣的问题：是什么使一个多粒子系统成为费米子系统？显然，仅有泡利不相容原理是不够的。一个费米系统不仅拥有满足泡利不相容原理的费米型希尔伯特空间，它的哈密顿量还具有一种非常特殊的性质。事实上，费米子系统是由高度非局域的哈密顿量
\begin{equation*}
H_f=\sum_{\langle\pmb{i}\pmb{j}\rangle}(\hat{t}_{\pmb{i}\pmb{j}}(\{\sigma_{\pmb{i}'}^{z}\})\sigma_{\pmb{i}}^{+}\sigma_{\pmb{j}}^{-}+\text{h.c.})
\end{equation*}
描述的，其中$\hat{t}_{\pmb{i}\pmb{j}}$是$\sigma_{\pmb{i}}^{z}$算符的函数，它包含许多个$\sigma_{\pmb{i}}^{z}$算符的乘积。对于一个有$N_s$个格点的$d$维格子，$\sigma_{\pmb{i}}^{z}$算符的数目是$N_s^{d-1}$的量级。如果不是大自然为我们提供了这样的非局域系统，任何头脑正常的物理学家都不会想去研究它们。既然这类非局域系统确实存在于自然界，我们就不得不研究它们。可是怎么研究呢？

这里我们极为幸运。自然界中的非局域费米子系统具有一些特殊的性质，使我们能够把它们简化。为了写下简化后的哈密顿量，我们先按某种方式给所有格点排序：$(\pmb{i}_1,\pmb{i}_2,\ldots,\pmb{i}_a,\ldots)$。然后我们引入另一类费米子算符：
\begin{equation*}
c_{\pmb{i}_a}=\sigma_{\pmb{i}_a}^{-}\prod_{b<a}\sigma_{\pmb{i}_b}^{z}\tag{4.1.1}
\end{equation*}
如果我们把格点排序得当，那么用$c_{\pmb{i}}$写出时，$H_f$将取如下较简单的形式：
\begin{equation*}
H_f=\sum_{\langle\pmb{i}\pmb{j}\rangle}(t_{\pmb{i}\pmb{j}}c_{\pmb{i}}^{\dagger}c_{\pmb{j}}+\text{h.c.})\tag{4.1.2}
\end{equation*}
其中$t_{\pmb{i}\pmb{j}}$不依赖于$\sigma_{\pmb{i}}^{z}=2c_{\pmb{i}}^{\dagger}c_{\pmb{i}}-1$。可以验证
\begin{align*}
&\{c_{\pmb{i}},c_{\pmb{j}}\}=\{c_{\pmb{i}}^{\dagger},c_{\pmb{j}}^{\dagger}\}=0\\
&\{c_{\pmb{i}},c_{\pmb{j}}^{\dagger}\}=\delta_{\pmb{i}\pmb{j}}\tag{4.1.3}
\end{align*}
其中$\{A,B\}\equiv AB+BA$称为反对易子（anti-commutator）。玻色子算符$\sigma^{x,y,z}$与费米子算符$c_{\pmb{i}}$之间的这一映射就是Jordan--Wigner变换（Jordan and Wigner, 1928）。我们注意到$c_{\pmb{i}}^{\dagger}c_{\pmb{i}}=\sigma_{\pmb{i}}^{+}\sigma_{\pmb{i}}^{-}$。这里$|0\rangle$和$|1\rangle$是$c_{\pmb{i}}^{\dagger}c_{\pmb{i}}$的两个本征态，本征值分别为0和1。因此，$c_{\pmb{i}}^{\dagger}c_{\pmb{i}}$是格点$\pmb{i}$上的费米子数算符。

通常，我们不去谈论费米子从何而来，而是直接把式(4.1.3)和式(4.1.2)当作自由费米子系统的定义。但如果我们真的追问费米子从何而来，并且认为玻色子比费米子更为基本，\footnote{玻色子系统的定义见10.1节。}那么上面的讨论表明，\emph{费米子是非局域激发}。事实上，自然界中的费米子的确表现得像非局域激发，因为费米子不能被单独产生。大自然似乎想要随时掌握我们宇宙中费米子的总数，以保证这个数是偶整数。局域激发不可能具有这样一种非局域的约束。在第10章中我们将看到，费米子可以解释为凝聚弦的端点，确实是非局域的。

\subsection{自由费米子系统的精确解}

\begin{keypoints}
\item 自由费米子系统的全部本征态和能量本征值都可以由费米子算符的反对易代数得到。
\end{keypoints}

在一个有$N_{\mathrm{site}}$个格点的格子上，费米子系统的希尔伯特空间有$2^{N_{\mathrm{site}}}$个态。哈密顿量(4.1.2)是一个$2^{N_{\mathrm{site}}}\times 2^{N_{\mathrm{site}}}$矩阵。求解该哈密顿量就相当于求出这么大一个矩阵的本征值和本征矢。

利用反对易代数(4.1.3)，可以把哈密顿量(4.1.2)严格解出。为简单起见，让我们假设系统具有平移对称性。此时$t_{\pmb{i}\pmb{j}}$只依赖于差$\pmb{i}-\pmb{j}$，即$t_{\pmb{i},\pmb{i}+\Delta\pmb{i}}=t_{\Delta\pmb{i}}$。引入$c_{\pmb{k}}=\sum_{\pmb{i}}N_s^{-1/2}\mathrm{e}^{-\mathrm{i}\pmb{k}\cdot\pmb{i}}c_{\pmb{i}}$，其中$N_s$是格点总数，我们发现
\begin{align*}
&H_f=\sum_{\pmb{k}}\epsilon_{\pmb{k}}c_{\pmb{k}}^{\dagger}c_{\pmb{k}},\qquad \epsilon_{\pmb{k}}=\sum_{\Delta\pmb{i}}t_{\Delta\pmb{i}}\mathrm{e}^{\mathrm{i}\pmb{k}\cdot\Delta\pmb{i}},\\
&\{c_{\pmb{k}},c_{\pmb{k}'}\}=\{c_{\pmb{k}}^{\dagger},c_{\pmb{k}'}^{\dagger}\}=0,\qquad \{c_{\pmb{k}},c_{\pmb{k}'}^{\dagger}\}=\delta_{\pmb{k}\pmb{k}'}.\tag{4.1.4}
\end{align*}
可以验证，满足$c_{\pmb{k}}|0\rangle=0$的态$|0\rangle$是$H_f$的零本征值本征态。利用反对易关系，我们发现态$|\{n_{\pmb{k}}\}\rangle\equiv\prod_{\pmb{k}}(c_{\pmb{k}}^{\dagger})^{n_{\pmb{k}}}|0\rangle$（$n_{\pmb{k}}=0,1$）是$c_{\pmb{k}}^{\dagger}c_{\pmb{k}}$的共同本征态，即$c_{\pmb{k}}^{\dagger}c_{\pmb{k}}|\{n_{\pmb{k}}\}\rangle=n_{\pmb{k}}|\{n_{\pmb{k}}\}\rangle$，而且它还是能量$E=\sum_{\pmb{k}}n_{\pmb{k}}\epsilon_{\pmb{k}}$的本征态。上面这种形式的本征态给出了$H_f$的全部本征态。这里的$n_{\pmb{k}}$是动量$\pmb{k}$上的费米子占据数。往$\pmb{k}_0$能级上加一个费米子，总能量就增加$\epsilon_{\pmb{k}_0}$，因此$\epsilon_{\pmb{k}_0}$可以解释为单粒子能量。

值得注意的是，我们通过代数途径就得到了$H_f$的全部本征态，而无需显式写出本征矢。对易的玻色代数和反对易的费米代数，是我们能够严格求解的仅有的几个代数系统中的两个。实际上，在相当长的时间里，它们曾是我们能在一维以外严格求解的仅有的两个代数系统。找到一维以外精确可解的代数系统，对我们理解相互作用系统极为重要。第10章将讨论最近发现的另外一些能够在一维以外求解的代数系统（Kitaev, 2003; Levin and Wen, 2003; Wen, 2003c）。这些代数系统引出了一维以外精确可解的相互作用模型。

如果计入化学势$\mu$，自由费米子系统的哈密顿量将变为
\begin{equation*}
H=\sum_{\pmb{k}}(\epsilon_{\pmb{k}}-\mu)c_{\pmb{k}}^{\dagger}c_{\pmb{k}}=\sum_{\pmb{k}}\xi_{\pmb{k}}c_{\pmb{k}}^{\dagger}c_{\pmb{k}},\qquad \xi_{\pmb{k}}\equiv\epsilon_{\pmb{k}}-\mu\tag{4.1.5}
\end{equation*}
$H$的基态由一个态$|\Psi_0\rangle$给出：所有$\xi_{\pmb{k}}<0$的$\pmb{k}$态都被填满，所有$\xi_{\pmb{k}}>0$的$\pmb{k}$态都空着。$|\Psi_0\rangle$由如下代数关系定义：
\begin{equation*}
\begin{aligned}
c_{\pmb{k}}|\Psi_0\rangle&=0,\qquad \text{当 }\xi_{\pmb{k}}>0\\
c_{\pmb{k}}^{\dagger}|\Psi_0\rangle&=0,\qquad \text{当 }\xi_{\pmb{k}}<0
\end{aligned}
\end{equation*}
占据数$n_{\pmb{k}}=c_{\pmb{k}}^{\dagger}c_{\pmb{k}}$在费米面处有一个跃变（见图4.1）。

%FIG{4.1}: {自由费米子系统的基态是一个费米海（Fermi sea）。这里的$E_F$称为费米能，$k_F$称为费米动量。占据数$n_k=n_F(\xi_k)$在费米动量$k_F$处不连续。}

\subsection*{问题 4.1.1}

\textbf{Jordan--Wigner变换}\quad 由式(4.1.1)证明式(4.1.3)。

\subsection*{问题 4.1.2}

\textbf{一维超导体的能谱}\quad 一维超导体由如下自由费米子哈密顿量描述：
\begin{equation*}
H=\sum_{i}(t\,c_{i}^{\dagger}c_{i+1}+\eta\,c_{i}^{\dagger}c_{i+1}^{\dagger}+\text{h.c.})\tag{4.1.6}
\end{equation*}
证明：在动量空间中，算符
\begin{equation*}
\lambda_{k}=u_{k}c_{k}+v_{k}c_{-k}^{\dagger}
\end{equation*}
满足费米子反对易关系$\{\lambda_{k},\lambda_{k'}\}=\{\lambda_{k}^{\dagger},\lambda_{k'}^{\dagger}\}=0$，并且当$|u_{k}|^{2}+|v_{k}|^{2}=1$时$\{\lambda_{k},\lambda_{k'}^{\dagger}\}=\delta_{kk'}$。证明：通过适当选取$u_{k}$和$v_{k}$，我们可以把$H$改写为
\begin{equation*}
H=E_{g}+\sum_{k}E_{k}\lambda_{k}^{\dagger}\lambda_{k}.
\end{equation*}
求出准粒子激发谱$E_{k}$和基态能量$E_{g}$。

\subsection*{问题 4.1.3}

\textbf{用Jordan--Wigner变换求解一维自旋-1/2系统（或硬核玻色子系统）}\quad 考虑一个具有如下哈密顿量的一维自旋-1/2系统（或硬核玻色子系统）：
\begin{equation*}
H=\sum_{i}\left(J_{x}\sigma_{i}^{x}\sigma_{i+1}^{x}+J_{y}\sigma_{i}^{y}\sigma_{i+1}^{y}+B\sigma_{i}^{z}\right)
\end{equation*}
\begin{enumerate}
\item 用Jordan--Wigner变换（对一维格子取自然排序），把上述相互作用的自旋-1/2（或硬核玻色子）系统映射为一个自由费米子系统。
\item 设$J_x=J_y=J$且$B\neq 0$。当$J$从$-\infty$变到$+\infty$时，系统经历几次相变。求出这些相变的临界$J$值。对每一个相和每一个临界点，画出总能量--动量空间中系统存在激发的区域。所有低能激发都应当包括在内，不论其动量多大；高能激发则不必包括。（提示：应当先思考并猜测一下，根据我们对一维相互作用玻色子系统的知识，谱应当是什么样子。）
\item 设$J_x=\alpha J_y=J$（$0<\alpha<1$）且$B\neq 0$。当$J$从$-\infty$变到$+\infty$时，系统经历几次相变。求出这些相变的临界$J$值。对每一个相和每一个临界点，画出总能量--动量空间中系统存在激发的区域。所有低能激发都应当包括在内，不论其动量多大。
\item 讨论为什么上述两种情形，即$J_x=J_y$与$J_x\neq J_y$，在定性上是不同的。
\end{enumerate}

\subsection{马约拉纳费米子}

自由费米子系统(4.1.2)是一个精确可解的多体系统，其希尔伯特空间的维数为$2^{N_{\mathrm{site}}}$（即每个格点两个态）。还有一个更小的精确可解系统，其希尔伯特空间中只有$2^{N_{\mathrm{site}}/2}$个态（假定$N_{\mathrm{site}}$为偶数）。这个系统称为马约拉纳费米子（Majorana fermion）系统，由如下哈密顿量描述：
\begin{equation*}
H=\sum_{i}(-\mathrm{i}t\,\lambda_{i}\lambda_{i+1}+\text{h.c.})\tag{4.1.7}
\end{equation*}
其中满足
\begin{equation*}
\{\lambda_{i},\lambda_{j}\}=\delta_{ij},\qquad (\lambda_{i})^{\dagger}=\lambda_{i}\tag{4.1.8}
\end{equation*}
的$\lambda_{i}$是马约拉纳费米子算符。马约拉纳费米子系统的希尔伯特空间构成上述反对易代数的一个表示。

为了构造希尔伯特空间，让我们考虑一个有$N_{\mathrm{site}}$个格点的一维周期系统，并假定$N_{\mathrm{site}}$为偶数，而且$\lambda_{i}$满足反周期边界条件$\lambda_{i+N_{\mathrm{site}}}=-\lambda_{i}$。在$k$空间中，哈密顿量(4.1.7)取如下形式：
\begin{equation*}
H=-\sum_{k}\mathrm{i}t\,\mathrm{e}^{\mathrm{i}k}\lambda(-k)\lambda(k)=\sum_{k>0}t\sin k\,\lambda^{\dagger}(k)\lambda(k)
\end{equation*}
其中$\lambda(k)=N_{\mathrm{site}}^{-1/2}\sum_{n}\mathrm{e}^{-\mathrm{i}kn}\lambda_{n}$，而$k=\pm\pi/N_{\mathrm{site}},\pm3\pi/N_{\mathrm{site}},\ldots,\pm(N_{\mathrm{site}}-1)\pi/N_{\mathrm{site}}$。可以验证
\begin{equation*}
\{\lambda^{\dagger}(k),\lambda(k')\}=\delta_{k-k'},\qquad \lambda^{\dagger}(k)=\lambda(-k)
\end{equation*}
我们看到，$k>0$的$\lambda(k)$满足复费米子的反对易代数。令$|0\rangle$为满足$\lambda(k)|0\rangle=0$（对$k>0$）的态。我们可以把$|0\rangle$看作没有费米子的态。令$|\{n_k\}\rangle=\prod_{k>0}(\lambda^{\dagger}(k))^{n_k}|0\rangle$，其中$n_k=0,1$。由于$\lambda^{\dagger}(k)\lambda(k)|\{n_k\}\rangle=n_k|\{n_k\}\rangle$，我们可以把$n_k$看作能级$k$上的占据数。注意，能级只有$N_{\mathrm{site}}/2$个，因为能级由正$k$标记（见图4.2）。因此，希尔伯特空间有$2^{N_{\mathrm{site}}/2}$个态。有趣的是，马约拉纳系统每个格点只有$\sqrt{2}$个态！态$|\{n_k\}\rangle$是一个能量本征态，能量为$\sum_{k>0}2t\sin(k)n_k$。

%FIG{4.2}: {马约拉纳费米子的单粒子能级仅对$k>0$存在。多体态由这些单粒子能级上的占据数$n_k$描述。}

对于一个具有两个马约拉纳费米子的系统
\begin{equation*}
H=\sum_{i}(-\mathrm{i}t\,\lambda_{1,i}\lambda_{1,i+1}-\mathrm{i}t\,\lambda_{2,i}\lambda_{2,i+1}+\text{h.c.}),
\end{equation*}
其多体希尔伯特空间有$2^{N_{\mathrm{site}}}$个态（即每个格点$\sqrt{2}\times\sqrt{2}=2$个态），与一个复费米子的希尔伯特空间相同。事实上，两个马约拉纳费米子的系统等价于一个复费米子的系统（见问题4.1.4）。

\subsection*{问题 4.1.4}

超导哈密顿量(4.1.6)也可以用马约拉纳费米子算符求解。马约拉纳费米子算符无非就是复费米子算符$c_{i}$的实部和虚部：
\begin{equation*}
\lambda_{1,i}=\frac{c_{i}+c_{i}^{\dagger}}{\sqrt{2}},\qquad \lambda_{2,i}=\frac{c_{i}-c_{i}^{\dagger}}{\mathrm{i}\sqrt{2}}
\end{equation*}
证明$\lambda_{ai}$满足马约拉纳费米子的反对易代数。用马约拉纳费米子求出所有能量本征态及其能量本征值，并把你的结果与问题4.1.2中得到的结果进行比较。

\subsection{跃迁算符的统计代数}

\begin{keypoints}
\item 全同粒子的统计性由其跃迁算符的代数决定。
\end{keypoints}

通常，玻色子系统被定义为由对易算符描述的系统，而费米子系统被定义为由反对易算符描述的系统。然而，这些定义太形式化了。为了从物理上理解玻色子系统与费米子系统的差别，我们考虑下面这个多体跃迁系统。希尔伯特空间由零粒子态$|0\rangle$、单粒子态$|\pmb{i}_1\rangle$、两粒子态$|\pmb{i}_1,\pmb{i}_2\rangle$等等构成，其中$\pmb{i}_n$标记格子中的格点。作为一个全同粒子系统，态$|\pmb{i}_1,\pmb{i}_2,\ldots\rangle$不依赖于指标$\pmb{i}_1,\pmb{i}_2,\ldots$的次序，例如$|\pmb{i}_1,\pmb{i}_2\rangle=|\pmb{i}_2,\pmb{i}_1\rangle$。没有双重占据的格点，并且我们约定当$\pmb{i}_m=\pmb{i}_n$时$|\pmb{i}_1,\pmb{i}_2,\ldots\rangle=0$。

跃迁算符$\hat{t}_{ij}$定义如下：当$\hat{t}_{ij}$作用在态$|\pmb{i}_1,\pmb{i}_2,\ldots\rangle$上时，如果格点$j$上有一个粒子而格点$i$上没有粒子，那么$\hat{t}_{ij}$就把$j$上的粒子搬到$i$，并在所得的态上乘一个复振幅$t(i,j;i_1,i_2,\ldots)$。注意，这个振幅可以依赖于所有粒子的位置，因而不一定是局域的。除此情形之外，跃迁算符$\hat{t}_{ij}$湮灭该态。我们系统的哈密顿量为
\begin{equation*}
H=\sum_{\langle ij\rangle}\hat{t}_{ij}
\end{equation*}
%FIG{4.3}: {（a）五次跃迁的第一种排序方式交换$i$、$j$处的两个粒子。（b）同样五次跃迁的第二种排序方式不交换这两个粒子。}
其中求和$\sum_{\langle ij\rangle}$遍及某一组配对$\langle ij\rangle$，例如最近邻配对。为了使上述哈密顿量描述一个局域系统，我们要求：当$i$、$j$、$k$、$l$互不相同时，
\begin{equation*}
[\hat{t}_{ij},\hat{t}_{kl}]=0
\end{equation*}
现在的问题是：上述跃迁哈密顿量描述的是硬核玻色子系统还是费米子系统？

多体跃迁系统是玻色子系统还是费米子系统（甚至是其他某种统计系统），与希尔伯特空间毫无关系。多体态由对称化的指标标记（例如$|\pmb{i}_1,\pmb{i}_2\rangle=|\pmb{i}_2,\pmb{i}_1\rangle$）这一事实，并不意味着多体系统就是玻色子系统，正如我们在4.1.1节中已经看到的。统计性由哈密顿量$H$决定。

显然，当跃迁振幅$t(i,j;i_1,i_2,\ldots)$只依赖于$i$和$j$，即$t(i,j;i_1,i_2,\ldots)=t(i,j)$时，多体跃迁哈密顿量描述的是硬核玻色子系统。问题是：在什么条件下多体跃迁哈密顿量描述费米子系统？

Levin和Wen（2003）解决了这个问题。他们发现，只要对任意三个跃迁算符$\hat{t}_{lj}$、$\hat{t}_{il}$、$\hat{t}_{lk}$（其中$i$、$j$、$k$、$l$互不相同），跃迁算符满足
\begin{equation*}
\hat{t}_{lk}\hat{t}_{il}\hat{t}_{lj}=-\hat{t}_{lj}\hat{t}_{il}\hat{t}_{lk}\tag{4.1.9}
\end{equation*}
多体跃迁哈密顿量就描述一个费米子系统。（注意该代数具有$\hat{t}_1\hat{t}_2\hat{t}_3=-\hat{t}_3\hat{t}_2\hat{t}_1$的结构。）

考虑态$|\pmb{i},\pmb{j},\ldots\rangle$，其中$i$、$j$处各有一个粒子，可能还有其他粒子在更远处。我们把一组五个跃迁算符$\{\hat{t}_{jl},\hat{t}_{lk},\hat{t}_{il},\hat{t}_{lj},\hat{t}_{ki}\}$按不同的次序作用到态$|\pmb{i},\pmb{j},\ldots\rangle$上（见图4.3）：
\begin{align*}
\hat{t}_{jl}\hat{t}_{lk}\hat{t}_{il}\hat{t}_{lj}\hat{t}_{ki}|\pmb{i},\pmb{j},\ldots\rangle&=C_1|\pmb{i},\pmb{j},\ldots\rangle\\
\hat{t}_{jl}\hat{t}_{lj}\hat{t}_{il}\hat{t}_{lk}\hat{t}_{ki}|\pmb{i},\pmb{j},\ldots\rangle&=C_2|\pmb{i},\pmb{j},\ldots\rangle
\end{align*}
其中我们假定格点$k$和$l$上没有粒子。我们注意到，五次跃迁之后我们回到原来的态$|\pmb{i},\pmb{j},\ldots\rangle$，只多出了相位$C_{1,2}$。然而从图4.3可以看到，五次跃迁的第一种排序方式（图4.3(a)）交换了$i$、$j$处的两个粒子，而第二种方式（图4.3(b)）不交换这两个粒子。既然两种跃迁方案使用的是同一组五次跃迁，$C_1$与$C_2$之间的差别就来自两个粒子的交换。因此，要使多体跃迁哈密顿量描述费米子系统，必须要求$C_1=-C_2$。注意到两种方案中第一次跃迁和最后一次跃迁是相同的，我们发现：只要跃迁算符满足式(4.1.9)，就有$C_1=-C_2$。事实上，如果我们不想使用反对易代数，式(4.1.9)可以作为费米统计的另一种定义。

我们想指出，在二维格子上，费米子跃迁代数可以推广为更一般的跃迁算符统计代数：
\begin{equation*}
\hat{t}_{ji}\hat{t}_{ik}\hat{t}_{li}=\mathrm{e}^{\mathrm{i}\theta}\hat{t}_{li}\hat{t}_{ik}\hat{t}_{ji}\tag{4.1.10}
\end{equation*}
在这种情形下，多体跃迁哈密顿量描述的是一个统计角为$\theta$的任意子（anyon）系统。

\subsection*{问题 4.1.5}

证明费米子跃迁算符$\hat{t}_{ij}=c_{i}^{\dagger}c_{j}$满足费米子跃迁代数。

\subsection*{问题 4.1.6}

我们已知弦算符（string operator）$\hat{W}_{ij}=\hat{t}_{il}\hat{t}_{lm}\ldots\hat{t}_{nj}$在$i$处产生一个粒子、在$j$处湮灭一个粒子，于是我们可能想把它写成$\hat{W}_{ij}=C_{i}^{+}C_{j}^{-}$。

\begin{enumerate}
\item[(a)] 证明：如果跃迁算符$\hat{t}_{ij}$满足费米子跃迁代数(4.1.9)，则弦算符满足$\hat{W}_{lk}\hat{W}_{il}\allowbreak\hat{W}_{lj}\allowbreak=-\hat{W}_{lj}\allowbreak\hat{W}_{il}\allowbreak\hat{W}_{lk}$。（可以假定这三条弦只在一点$l$相交。）
\item[(b)] 证明$C_{i}^{\pm}$与$C_{j}^{\pm}$不能对易，即使$i$与$j$相距很远。但是，如果$C_{i}^{\pm}$与$C_{j}^{\pm}$反对易（当$i\neq j$时），则费米子跃迁代数(4.1.9)可以得到满足。
\end{enumerate}

\section{自由费米子格林函数}

\begin{keypoints}
\item 格林函数在很长时间上的代数衰减与跨费米面的无能隙激发有关；长距离上的代数衰减与动量空间中的尖锐特征（不连续性）有关。
\item 电子格林函数可以在隧穿实验中测量。
\item 电子格林函数的谱函数。
\end{keypoints}

在接下来的几节中，我们将讨论自由费米子系统的单体和两体关联函数。单体关联函数对理解隧穿和光电子发射（photoemission）实验十分重要。两体关联函数则更为重要，因为它们与各种各样的输运、散射以及线性响应实验都有关系。在这两节中我会给出一些计算细节。这些讨论的目的主要是引进数学形式。我们列出了$d=1,2,3$维的许多显式结果，希望这些结果既可以用作参考，也可以作为各种关联的具体例子。

\subsection{编时关联函数}

\begin{keypoints}
\item 在编时平均之下，费米子算符表现得像反对易的数。
\end{keypoints}

考虑由式(4.1.5)描述的自由费米子系统。在海森堡绘景中，含时的费米子算符为
\begin{equation*}
c_{\pmb{k}}(t)=U^{\dagger}(t,-\infty)c_{\pmb{k}}U(t,-\infty),\qquad U(t_1,t_2)=\mathrm{e}^{-\mathrm{i}\int_{t_2}^{t_1}\mathrm{d}t\, H}
\end{equation*}
零温和有限温下的费米子传播子（即格林函数）定义为编时平均
\begin{align*}
\mathrm{i}G(t_1-t_2,\pmb{k}_1)\delta_{\pmb{k}_1,\pmb{k}_2}&=\langle 0|T(c_{\pmb{k}_1}(t_1)c_{\pmb{k}_2}^{\dagger}(t_2))|0\rangle\tag{4.2.1}\\
&=\begin{cases}
+\langle 0|c_{\pmb{k}_1}(t_1)c_{\pmb{k}_2}^{\dagger}(t_2)|0\rangle, & \text{当 } t_1>t_2\\
-\langle 0|c_{\pmb{k}_2}^{\dagger}(t_2)c_{\pmb{k}_1}(t_1)|0\rangle, & \text{当 } t_1<t_2
\end{cases}\qquad \text{当 } T=0
\end{align*}
\begin{align*}
\mathrm{i}G^{\beta}(t_1-t_2,\pmb{k}_1)\delta_{\pmb{k}_1,\pmb{k}_2}&=\langle T(c_{\pmb{k}_1}(t_1)c_{\pmb{k}_2}^{\dagger}(t_2))\rangle\\
&\equiv\frac{\mathrm{Tr}[T(c_{\pmb{k}_1}(t_1)c_{\pmb{k}_2}^{\dagger}(t_2))\mathrm{e}^{-\beta H}]}{Z}\tag{4.2.2}\\
&=\begin{cases}
+\mathrm{Tr}[c_{\pmb{k}_1}(t_1)c_{\pmb{k}_2}^{\dagger}(t_2)\mathrm{e}^{-\beta H}], & \text{当 } t_1>t_2\\
-\mathrm{Tr}[c_{\pmb{k}_2}^{\dagger}(t_2)c_{\pmb{k}_1}(t_1)\mathrm{e}^{-\beta H}], & \text{当 } t_1<t_2
\end{cases}\qquad \text{当 } T>0
\end{align*}
注意上述定义中的负号；并且按定义
\begin{equation*}
\left\langle T\left(c^{\dagger}(t)c(t')\right)\right\rangle=-\left\langle T\left(c(t')c^{\dagger}(t)\right)\right\rangle
\end{equation*}
已知$\pmb{k}$态上的费米子占据数
\begin{equation*}
n_F(\xi_{\pmb{k}})=\langle c_{\pmb{k}}^{\dagger}c_{\pmb{k}}\rangle=\frac{1}{\mathrm{e}^{\beta\xi_{\pmb{k}}}+1}
\end{equation*}
我们就可以算出动量空间中的零温格林函数$G$和有限温格林函数$G^{\beta}$：
\begin{align*}
\mathrm{i}G(t,\pmb{k})&=+\Theta(t)\Theta(\xi_{\pmb{k}})\mathrm{e}^{-\mathrm{i}t\xi_{\pmb{k}}}-\Theta(-t)\Theta(-\xi_{\pmb{k}})\mathrm{e}^{-\mathrm{i}(-t)(-\xi_{\pmb{k}})}\big|_{T=0}\\
\mathrm{i}G^{\beta}(t,\pmb{k})&=+\Theta(t)(1-n_F(\xi_{\pmb{k}}))\mathrm{e}^{-\mathrm{i}t\xi_{\pmb{k}}}-\Theta(-t)n_F(\xi_{\pmb{k}})\mathrm{e}^{-\mathrm{i}t\xi_{\pmb{k}}}\big|_{T>0}\tag{4.2.3}
\end{align*}
在$\omega$--$\pmb{k}$空间中，$G(\omega,\pmb{k})=\int\mathrm{d}t\, G(t,\pmb{k})\mathrm{e}^{\mathrm{i}\omega t}$与$G^{\beta}(\omega,\pmb{k})=\int\mathrm{d}t\, G^{\beta}(t,\pmb{k})\mathrm{e}^{\mathrm{i}\omega t}$具有如下更简单的形式：
\begin{equation*}
G(\omega,\pmb{k})=\frac{1}{\omega-\xi_{\pmb{k}}+\mathrm{i}0^{+}\,\mathrm{sgn}(\omega)}\qquad \text{当 } T=0\tag{4.2.4}
\end{equation*}
\begin{equation*}
G^{\beta}(\omega,\pmb{k})=\frac{1-n_F(\xi_{\pmb{k}})}{\omega-\xi_{\pmb{k}}+\mathrm{i}0^{+}}+\frac{n_F(\xi_{\pmb{k}})}{\omega-\xi_{\pmb{k}}-\mathrm{i}0^{+}}\qquad \text{当 } T>0\tag{4.2.5}
\end{equation*}
现在让我们来考虑有限温下的虚时格林函数。此时，海森堡绘景中的含时算符为
\begin{equation*}
c_{\pmb{k}}(\tau)=\mathrm{e}^{H\tau}c_{\pmb{k}}\mathrm{e}^{-H\tau}
\end{equation*}
由定义（$0<\tau_1,\tau_2<\beta$）
\begin{equation*}
\mathcal{G}^{\beta}(\tau_1,\tau_2,\pmb{k}_1)\delta_{\pmb{k}_1,\pmb{k}_2}
=\begin{cases}
+\dfrac{\mathrm{Tr}(\mathrm{e}^{-\beta H}c_{\pmb{k}_1}(\tau_1)c_{\pmb{k}_2}^{\dagger}(\tau_2))}{\mathrm{Tr}(\mathrm{e}^{-\beta H})}, & \text{当 } \tau_1>\tau_2\\[10pt]
-\dfrac{\mathrm{Tr}(\mathrm{e}^{-\beta H}c_{\pmb{k}_2}(\tau_2)c_{\pmb{k}_1}^{\dagger}(\tau_1))}{\mathrm{Tr}(\mathrm{e}^{-\beta H})}, & \text{当 } \tau_1<\tau_2
\end{cases}\tag{4.2.6}
\end{equation*}
可以证明（见问题4.2.1）
\begin{align*}
&\mathcal{G}^{\beta}(\tau_1,\tau_2,\pmb{k})=\mathcal{G}^{\beta}(\tau_1-\tau_2,0,\pmb{k})\equiv\mathcal{G}^{\beta}(\tau_1-\tau_2,\pmb{k})\\
&\mathcal{G}^{\beta}(\tau,\pmb{k})=-\mathcal{G}^{\beta}(\tau+\beta,\pmb{k})\tag{4.2.7}
\end{align*}
因此，费米子格林函数在紧致化的虚时方向上是反周期的。对于自由费米子，我们有
\begin{equation*}
\mathcal{G}^{\beta}(\tau,\pmb{k})=+\Theta(\tau)(1-n_F(\xi_{\pmb{k}}))\mathrm{e}^{-\tau\xi_{\pmb{k}}}-\Theta(-\tau)n_F(\xi_{\pmb{k}})\mathrm{e}^{-\tau\xi_{\pmb{k}}}\tag{4.2.8}
\end{equation*}
在$\omega$--$\pmb{k}$空间中，有
\begin{equation*}
\mathcal{G}^{\beta}(\omega_{\gamma},\pmb{k})\equiv\int_{0}^{\beta}\mathrm{d}\tau\, \mathcal{G}^{\beta}(\tau,\pmb{k})\mathrm{e}^{\mathrm{i}\omega_{\gamma}\tau}=\frac{1}{-\mathrm{i}\omega_{\gamma}+\xi_{\pmb{k}}}\tag{4.2.9}
\end{equation*}
其中$\omega_{\gamma}=2\pi\gamma T$，这里$\gamma=\frac{1}{2}+$整数。

$\omega$--$\pmb{k}$空间中的编时格林函数可以写成$\omega$--$\pmb{k}$空间中费米子算符的\emph{编时}平均。对于虚时，我们引入
\begin{equation*}
c_{\omega_{\gamma}\pmb{k}}=\int_{0}^{\beta}\mathrm{d}\tau\, c_{\pmb{k}}(\tau)\beta^{-1/2}\mathrm{e}^{\mathrm{i}\omega_{\gamma}\tau}
\end{equation*}
经过一点运算后，我们发现
\begin{align*}
\langle T(c_{\omega_{\gamma}\pmb{k}}c_{\omega_{\gamma}'\pmb{k}}^{\dagger})\rangle&\equiv\beta^{-1}\int_{0}^{\beta}\mathrm{d}\tau\,\mathrm{d}\tau'\,\mathrm{e}^{\mathrm{i}\omega_{\gamma}\tau-\mathrm{i}\omega_{\gamma}'\tau'}\left\langle T\left(c_{\pmb{k}}(\tau)c_{\pmb{k}}^{\dagger}(\tau')\right)\right\rangle\\
&=\mathcal{G}^{\beta}(\omega_{\gamma},\pmb{k})\delta_{\omega_{\gamma}-\omega_{\gamma}'}\tag{4.2.10}
\end{align*}

%-----------------------------------------------------------------------------
% 新增术语（ch4_a 块新增，供汇总入术语表）：
%   hard-core boson system                    硬核玻色子系统
%   Jordan--Wigner transformation             Jordan--Wigner 变换
%   anti-commutator                           反对易子
%   Fermi sea                                 费米海
%   Fermi momentum                            费米动量
%   Fermi statistics                          费米统计
%   fermion number operator                   费米子数算符
%   single-particle energy                    单粒子能量
%   single-particle level                     单粒子能级
%   hopping operator / hopping algebra        跃迁算符 / 跃迁代数
%   string operator                           弦算符
%   statistical angle                         统计角
%   imaginary-time Green's function           虚时格林函数
%   anti-periodic boundary condition          反周期边界条件
%   compactified imaginary-time direction     紧致化的虚时方向
%   photoemission experiment                  光电子发射实验
%   time-ordered average                      编时平均
%   fermion propagator                        费米子传播子
%   quasiparticle excitation spectrum         准粒子激发谱
%   ground-state energy                       基态能量
%   tunneling Hamiltonian                     隧穿哈密顿量
%   total energy--momentum space              总能量--动量空间
